\chapter{General Topology}\label{ch:b3:topology}

The Year 2 volume did analysis in metric spaces: distances,
balls, sequences. But the fundamental notions --- \omterm{def:b3:topology:continuity}{continuity},
compactness, \omterm{def:b3:topology:connected}{connectedness} --- never mention the numerical value
of a distance, only the family of \emph{\omterm{def:b3:topology:topology}{open sets}} it generates.
This chapter takes that family as the primitive object. The gain
is not generality for its own sake: quotient constructions (the
circle as $\R/\Z$, projective spaces), products, and the
weak-type topologies of functional analysis simply are not
metric-first objects. We rebuild \omterm{def:b3:topology:continuity}{continuity}, then treat
compactness by open covers (proving it equivalent to Year 2's
sequential definition in metric spaces), and \omterm{def:b3:topology:connected}{connectedness} ---
closing with the theorem that a \omterm{def:b3:topology:continuity}{continuous} bijection cannot
identify $\R$ with $\R^2$: \omterm{def:b3:topology:topology}{topology} can distinguish dimensions.

\section{Topologies, open sets, continuity}

\begin{definition}\label{def:b3:topology:topology}
A \emph{topology}\index{topology} on a set $X$ is a family
$\mathcal T$ of subsets of $X$ --- called \emph{open
sets}\index{open set} --- such that: $\varnothing, X \in
\mathcal T$; any union of open sets is open; any \emph{finite}
intersection of open sets is open. The pair $(X, \mathcal T)$ is
a \emph{topological space}. Complements of open sets are
\emph{closed}. A \emph{neighborhood}\index{neighborhood} of $x$
is a set containing an open set containing $x$.
\end{definition}

\begin{example}\label{ex:b3:topology:examples}
(a) A metric space, with ``open'' as in Year 2 (unions of open
balls): the \emph{metric \omterm{def:b3:topology:topology}{topology}}; different metrics can give
the same \omterm{def:b3:topology:topology}{topology} (equivalent metrics). A space whose \omterm{def:b3:topology:topology}{topology}
arises from some metric is \emph{metrizable}.
(b) The \emph{discrete} \omterm{def:b3:topology:topology}{topology} (all subsets) and the
\emph{indiscrete} \omterm{def:b3:topology:topology}{topology} $\{\varnothing, X\}$.
(c) The \emph{cofinite} \omterm{def:b3:topology:topology}{topology} on an infinite set: open $=$
empty or cofinite. Not metrizable, as we shall see
(\cref{exo:b3:topology:3}).
(d) On $\R$, the usual \omterm{def:b3:topology:topology}{topology}; on $\bar\R = \R \cup
\{\pm\infty\}$, the order \omterm{def:b3:topology:topology}{topology} generated by rays --- making
``$x_n \to +\infty$'' an instance of plain convergence.
\end{example}

\begin{definition}\label{def:b3:topology:interior}
For $A \subseteq X$: the \emph{interior} $\mathring A$ is the
largest \omterm{def:b3:topology:topology}{open set} inside $A$ (union of all of them); the
\emph{closure}\index{closure} $\bar A$ the smallest closed set
containing $A$; the \emph{boundary} $\partial A = \bar A
\setminus \mathring A$. $A$ is \emph{dense}\index{dense set} if
$\bar A = X$. One has $x \in \bar A$ iff every \omterm{def:b3:topology:topology}{neighborhood} of
$x$ meets $A$ (if some \omterm{def:b3:topology:topology}{neighborhood} misses $A$, its open core's
complement is a smaller closed set around $A$; conversely).
\end{definition}

\begin{definition}\label{def:b3:topology:basis}
A \emph{basis}\index{basis of a topology} of $\mathcal T$ is a
family $\mathcal B \subseteq \mathcal T$ such that every \omterm{def:b3:topology:topology}{open
set} is a union of members of $\mathcal B$ (e.g.\ open balls in a
metric space; open intervals in $\R$). A family $\mathcal B$ of
subsets of $X$ is a basis of \emph{some} \omterm{def:b3:topology:topology}{topology} iff it covers
$X$ and, for $B_1, B_2 \in \mathcal B$ and $x \in B_1 \cap B_2$,
some $B_3 \in \mathcal B$ has $x \in B_3 \subseteq B_1\cap B_2$
--- then ``unions of members'' is a \omterm{def:b3:topology:topology}{topology}, the \omterm{def:b3:topology:topology}{topology}
\emph{generated} by $\mathcal B$.
\end{definition}

\begin{definition}\label{def:b3:topology:continuity}
$f \colon X \to Y$ is \emph{continuous}\index{continuous map} if
$f^{-1}(V)$ is open for every open $V \subseteq Y$ ---
equivalently, preimages of closed sets are closed; equivalently,
for every $x$ and \omterm{def:b3:topology:topology}{neighborhood} $V$ of $f(x)$, $f^{-1}(V)$ is a
\omterm{def:b3:topology:topology}{neighborhood} of $x$ (continuity \emph{at} each $x$). It
suffices to check on a \omterm{def:b3:topology:basis}{basis} of $Y$. Compositions of continuous
maps are continuous. A \emph{homeomorphism}\index{homeomorphism}
is a continuous bijection with continuous inverse; \omterm{def:b3:topology:topology}{topology}
studies the properties preserved by homeomorphisms.
\end{definition}

\begin{proposition}[Continuity vs closure; gluing]
\label{prop:b3:topology:gluing}
(a) $f$ is \omterm{def:b3:topology:continuity}{continuous} iff $f(\bar A) \subseteq \overline{f(A)}$
for all $A \subseteq X$.
(b) If $X = F_1 \cup F_2$ with $F_1, F_2$ closed and $f\colon X
\to Y$ restricts \omterm{def:b3:topology:continuity}{continuously} to each $F_i$, then $f$ is
\omterm{def:b3:topology:continuity}{continuous}.
\end{proposition}

\begin{proof}
(a) If $f$ \omterm{def:b3:topology:continuity}{continuous}: $f^{-1}(\overline{f(A)})$ is closed and
contains $A$, hence contains $\bar A$. Conversely apply the
criterion to $A = f^{-1}(F)$, $F$ closed: $f(\bar A) \subseteq
\overline{f(f^{-1}(F))} \subseteq \bar F = F$, so $\bar A
\subseteq f^{-1}(F) = A$: preimages of closed sets are closed.
(b) For $F \subseteq Y$ closed: $f^{-1}(F) =
(f\restriction_{F_1})^{-1}(F) \cup (f\restriction_{F_2})^{-1}
(F)$, a union of two sets closed in $F_1$, resp.\ $F_2$, hence
closed in $X$ ($F_i$ closed: closed-in-closed is closed).
\end{proof}

\begin{definition}\label{def:b3:topology:hausdorff}
$X$ is \emph{Hausdorff}\index{Hausdorff space} (or
\emph{separated}) if any two distinct points have disjoint
\omterm{def:b3:topology:topology}{neighborhoods}. Metric spaces are Hausdorff (balls of radius
$d(x,y)/2$). A sequence $(x_n)$ \emph{converges} to $x$ if every
\omterm{def:b3:topology:topology}{neighborhood} of $x$ contains all but finitely many $x_n$; in a
Hausdorff space, limits are unique (two limits would have
disjoint \omterm{def:b3:topology:topology}{neighborhoods} each containing a tail). In a Hausdorff
space, points --- hence finite sets --- are closed.
\end{definition}

\begin{remark}\label{rem:b3:topology:sequences}
In metric spaces, sequences detect everything: $x \in \bar A$
iff some sequence of $A$ converges to $x$ (take $x_n \in A \cap
B(x, 1/n)$), and $f$ is \omterm{def:b3:topology:continuity}{continuous} iff it is sequentially
\omterm{def:b3:topology:continuity}{continuous}. In general spaces both equivalences fail; the
correct sequence-substitutes (filters, nets) belong to a more
advanced course. We will state sequence-based results in metric
spaces and cover-based results in general --- and prove them
equivalent where they are.
\end{remark}

\section{Subspaces, products, quotients}

\begin{definition}\label{def:b3:topology:constructions}
Three ways to make new spaces from old:\index{product
topology}\index{quotient topology}
\begin{enumerate}
  \item \emph{Subspace}: on $A \subseteq X$, the \omterm{def:b3:topology:topology}{open sets} are
        the $U \cap A$, $U$ open in $X$ --- the coarsest
        \omterm{def:b3:topology:topology}{topology} making the inclusion \omterm{def:b3:topology:continuity}{continuous}.
  \item \emph{Product}: on $X \times Y$ (and finite products),
        the \omterm{def:b3:topology:topology}{topology} with \omterm{def:b3:topology:basis}{basis} the \emph{open boxes} $U \times
        V$; on an infinite product $\prod_i X_i$, the \omterm{def:b3:topology:basis}{basis}
        consists of boxes $\prod U_i$ with $U_i = X_i$ for
        \emph{all but finitely many} $i$ --- the coarsest
        \omterm{def:b3:topology:topology}{topology} making every projection \omterm{def:b3:topology:continuity}{continuous}.
  \item \emph{Quotient}: if $\sim$ is an equivalence on $X$ and
        $\pi\colon X \to X/{\sim}$ the projection, declare $V
        \subseteq X/{\sim}$ open iff $\pi^{-1}(V)$ is open ---
        the finest \omterm{def:b3:topology:topology}{topology} making $\pi$ \omterm{def:b3:topology:continuity}{continuous}.
\end{enumerate}
\end{definition}

\begin{proposition}[Universal properties]
\label{prop:b3:topology:universal}
(a) $f \colon Z \to X \times Y$ is \omterm{def:b3:topology:continuity}{continuous} iff both
components $\pi_X \circ f$, $\pi_Y \circ f$ are (same for
arbitrary products).
(b) $g \colon X/{\sim} \to Z$ is \omterm{def:b3:topology:continuity}{continuous} iff $g \circ \pi
\colon X \to Z$ is.
\end{proposition}

\begin{proof}
(a) Necessity: compositions. Sufficiency: it is enough to check
preimages of \omterm{def:b3:topology:basis}{basis} boxes: $f^{-1}(U \times V) = (\pi_Xf)^{-1}(U)
\cap (\pi_Yf)^{-1}(V)$, open (finitely many factors $\neq X_i$
in the infinite case). (b) Necessity: composition. Sufficiency:
for $W \subseteq Z$ open, $\pi^{-1}(g^{-1}(W)) = (g\pi)^{-1}(W)$
is open, which by definition of the \omterm{def:b3:topology:constructions}{quotient topology} means
$g^{-1}(W)$ is open.
\end{proof}

\begin{example}\label{ex:b3:topology:circle}
The quotient $\R/\Z$ (identify $x$ and $x + n$) is \omterm{def:b3:topology:continuity}{homeomorphic}
to the circle $S^1 = \{z \in \C : \abs z = 1\}$: the map $x
\mapsto \eu^{2\iu\pi x}$ passes to a \omterm{def:b3:topology:continuity}{continuous} bijection
$\R/\Z \to S^1$ (\cref{prop:b3:topology:universal}(b)); its
inverse is \omterm{def:b3:topology:continuity}{continuous} by the compactness argument of
\cref{cor:b3:topology:compacthomeo} below
(\cref{exo:b3:topology:5} details everything, including why
$\R/\Z$ is \omterm{def:b3:topology:hausdorff}{Hausdorff} and \omterm{def:b3:topology:compact}{compact}). Likewise $[0,1]$ with
endpoints glued is $S^1$, the square with opposite sides glued
is the torus, and gluing is finally a theorem, not a picture.
\end{example}

\section{Compactness}

\begin{definition}\label{def:b3:topology:compact}
An \emph{open cover} of $X$ is a family $(U_i)_{i\in I}$ of \omterm{def:b3:topology:topology}{open
sets} with $\bigcup U_i = X$. $X$ is
\emph{compact}\index{compact space} if it is \omterm{def:b3:topology:hausdorff}{Hausdorff} and every
open cover admits a \emph{finite subcover}. Equivalently (taking
complements): every family of closed sets with the \emph{finite
intersection property} (all finite subfamilies have nonempty
intersection) has nonempty total intersection.
\end{definition}

\begin{theorem}[First properties]\label{thm:b3:topology:compactprops}
Let $X$ be \omterm{def:b3:topology:compact}{compact}.
\begin{enumerate}
  \item A closed subset of $X$ is \omterm{def:b3:topology:compact}{compact}; a \omterm{def:b3:topology:compact}{compact} subset of a
        \omterm{def:b3:topology:hausdorff}{Hausdorff space} is closed.
  \item A \omterm{def:b3:topology:continuity}{continuous} image of a \omterm{def:b3:topology:compact}{compact space} in a \omterm{def:b3:topology:hausdorff}{Hausdorff
        space} is \omterm{def:b3:topology:compact}{compact}. In particular a \omterm{def:b3:topology:continuity}{continuous} $f\colon X
        \to \R$ is bounded and attains its bounds.
  \item A decreasing sequence of nonempty closed subsets of $X$
        has nonempty intersection.
\end{enumerate}
\end{theorem}

\begin{proof}
(1) Let $F \subseteq X$ closed and $(U_i)$ an open cover of $F$
(by opens of $X$): adding $X \setminus F$ gives an open cover of
$X$; a finite subcover, minus $X \setminus F$, covers $F$.
Subspace \omterm{def:b3:topology:hausdorff}{Hausdorff} is clear. Conversely let $K \subseteq Y$
\omterm{def:b3:topology:compact}{compact}, $Y$ \omterm{def:b3:topology:hausdorff}{Hausdorff}, and $y \notin K$: for each $x \in K$
pick disjoint open $U_x \ni x$, $V_x \ni y$; finitely many $U_x$
cover $K$, and the intersection of the corresponding $V_x$ is a
\omterm{def:b3:topology:topology}{neighborhood} of $y$ disjoint from $K$'s cover, hence from $K$:
the complement of $K$ is open.

(2) If $(V_j)$ covers $f(X)$, then $(f^{-1}(V_j))$ covers $X$; a
finite subcover downstairs comes from the finite subcover
upstairs. The image is \omterm{def:b3:topology:hausdorff}{Hausdorff} as a subspace. For $f$ real:
$f(X)$ is \omterm{def:b3:topology:compact}{compact} in $\R$, hence closed and bounded (cover by
$(-n, n)$ for boundedness; closed by (1)), and a closed bounded
set contains its supremum.

(3) If $\bigcap F_n = \varnothing$, the opens $X \setminus F_n$
cover $X$; finitely many suffice, so some $F_{n_0} =
\varnothing$ (the sequence decreases) --- contradiction.
\end{proof}

\begin{corollary}\label{cor:b3:topology:compacthomeo}
A \omterm{def:b3:topology:continuity}{continuous} \emph{bijection} from a \omterm{def:b3:topology:compact}{compact space} to a
\omterm{def:b3:topology:hausdorff}{Hausdorff space} is a \omterm{def:b3:topology:continuity}{homeomorphism}.
\end{corollary}

\begin{proof}
The inverse is \omterm{def:b3:topology:continuity}{continuous} iff direct images of closed sets are
closed; a closed $F$ is \omterm{def:b3:topology:compact}{compact}
(\cref{thm:b3:topology:compactprops}(1)), its image is \omterm{def:b3:topology:compact}{compact}
(2), hence closed (1) in the \omterm{def:b3:topology:hausdorff}{Hausdorff} target.
\end{proof}

\begin{theorem}[Finite products]\label{thm:b3:topology:tychonoff}
A finite product of \omterm{def:b3:topology:compact}{compact spaces} is \omterm{def:b3:topology:compact}{compact}.\index{Tychonoff's
theorem (finite case)}
\end{theorem}

\begin{proof}
It suffices to treat $X \times Y$. \omterm{def:b3:topology:hausdorff}{Hausdorff} is inherited
(separate in one coordinate). Let $(W_i)$ be an open cover of $X
\times Y$; we may assume the $W_i$ are \omterm{def:b3:topology:basis}{basis} boxes $U_i \times
V_i$ (refine: each point sits in a box inside some $W_i$; a
finite subcover of boxes yields one of $W_i$'s). Fix $x \in X$:
the \emph{slice} $\{x\}\times Y \cong Y$ is \omterm{def:b3:topology:compact}{compact}, so finitely
many boxes $U_1\times V_1, \dots, U_k\times V_k$ cover it, with
$x \in U_j$ for all $j$; then $U^x = \bigcap_{j \leq k} U_j$ is
an open \omterm{def:b3:topology:topology}{neighborhood} of $x$ with $U^x \times Y$ covered by
finitely many boxes (the \emph{tube lemma}\index{tube lemma}:
for $(x', y) \in U^x \times Y$, $y \in V_j$
for some $j$ since the boxes covered $\{x\}\times Y$ at level
$y$, and $x' \in U^x \subseteq U_j$, so $(x', y) \in U_j
\times V_j$). Now finitely many $U^x$
cover the \omterm{def:b3:topology:compact}{compact} $X$; the corresponding finite collections of
boxes cover $X \times Y$.
\end{proof}

\begin{theorem}[Compactness in metric spaces]
\label{thm:b3:topology:metriccompact}
For a metric space $(X, d)$, the following are
equivalent:\index{sequential compactness}
\begin{enumerate}
  \item $X$ is \omterm{def:b3:topology:compact}{compact} (Borel--Lebesgue);
  \item every sequence in $X$ has a convergent subsequence
        (\omterm{thm:b3:topology:metriccompact}{sequential compactness} --- Year 2's definition);
  \item $X$ is complete and \emph{totally bounded}: for every
        $\varepsilon > 0$, finitely many balls of radius
        $\varepsilon$ cover $X$.
\end{enumerate}
\end{theorem}

\begin{proof}
(1)$\Rightarrow$(2): Let $(x_n)$ have no convergent
subsequence. Then each $x \in X$ has an open ball $B_x$
containing $x_n$ for only finitely many $n$ (otherwise a
subsequence converges to $x$: take radii $1/k$). Finitely many
$B_x$ cover $X$, so only finitely many indices $n$ exist:
absurd.

(2)$\Rightarrow$(3): Completeness: a Cauchy sequence with a
convergent subsequence converges (Year 2). Total boundedness: if
some $\varepsilon$ admits no finite cover, choose inductively
$x_{n+1}$ outside $\bigcup_{k\leq n} B(x_k, \varepsilon)$: the
sequence has $d(x_m, x_n) \geq \varepsilon$ for $m \neq n$, no
Cauchy subsequence, no convergent one.

(3)$\Rightarrow$(1): First, (3) implies (2): given $(x_n)$,
cover $X$ by finitely many balls of radius $1$: one, $B_1$,
contains a subsequence; cover $X$ by balls of radius $1/2$: one
contains a further subsequence; iterate and diagonalize: the
diagonal subsequence is Cauchy (two terms beyond stage $k$ lie
in a common ball of radius $2^{-k}$, up to the usual $2\cdot$),
hence converges. Now let $(U_i)$ be an open cover and suppose no
finite subcover. \emph{Lebesgue number argument}: for each $n$,
some ball $B(y_n, 2^{-n})$ is not covered by finitely many
$U_i$ --- indeed, cover $X$ by finitely many balls of radius
$2^{-n}$; if each were finitely covered, so would $X$ be. By
(2), a subsequence $y_{n_k} \to y$; pick $i$ with $y \in U_i$
and $r > 0$ with $B(y, r) \subseteq U_i$. For large $k$,
$B(y_{n_k}, 2^{-n_k}) \subseteq B(y, r) \subseteq U_i$: covered
by \emph{one} $U_i$ --- contradiction.
\end{proof}

\begin{corollary}[Heine--Borel; Heine]
\label{cor:b3:topology:heineborel}
(a) A subset of $\R^n$ is \omterm{def:b3:topology:compact}{compact} iff it is closed and bounded.
(b) A \omterm{def:b3:topology:continuity}{continuous map} from a \omterm{def:b3:topology:compact}{compact} metric space to a metric
space is uniformly \omterm{def:b3:topology:continuity}{continuous}.\index{Heine--Borel
theorem}\index{Heine's theorem}
\end{corollary}

\begin{proof}
(a) Closed and bounded $\Rightarrow$ contained in a cube
$[-M,M]^n$, which is \omterm{def:b3:topology:compact}{compact}: $[-M, M]$ is (sequentially, by
Bolzano--Weierstrass --- or directly by dichotomy for covers),
and \cref{thm:b3:topology:tychonoff} handles the product; then
apply \cref{thm:b3:topology:compactprops}(1). Conversely a
\omterm{def:b3:topology:compact}{compact} subset is closed
(\cref{thm:b3:topology:compactprops}(1)) and bounded (cover by
concentric balls).

(b) Let $f \colon X \to Y$, $\varepsilon > 0$. The balls
$B\bigl(x, \delta_x\bigr)$ with $f\bigl(B(x,
2\delta_x)\bigr)\subseteq B\bigl(f(x), \varepsilon/2\bigr)$
cover $X$; extract a finite subcover $B(x_j, \delta_{x_j})$ and
set $\delta = \min_j \delta_{x_j}$. If $d(x, x') < \delta$: $x
\in B(x_j, \delta_{x_j})$ for some $j$, and then both $x, x'
\in B(x_j, 2\delta_{x_j})$, so $d(f(x), f(x')) < \varepsilon$.
\end{proof}

\begin{definition}\label{def:b3:topology:locallycompact}
$X$ is \emph{locally compact}\index{locally compact space} if it
is \omterm{def:b3:topology:hausdorff}{Hausdorff} and every point has a \omterm{def:b3:topology:compact}{compact} \omterm{def:b3:topology:topology}{neighborhood} ($\R^n$;
open subsets of $\R^n$; discrete spaces --- but not $\Q$, see
\cref{exo:b3:topology:9}). Every locally \omterm{def:b3:topology:compact}{compact space} embeds in
a \omterm{def:b3:topology:compact}{compact} one: the \emph{one-point
compactification}\index{one-point compactification} $\hat X = X
\cup \{\infty\}$, whose opens are those of $X$ together with the
complements (in $\hat X$) of \omterm{def:b3:topology:compact}{compact} subsets of $X$. One checks
the axioms directly; $\hat X$ is \omterm{def:b3:topology:compact}{compact} (a cover has a member
containing $\infty$, whose complement is \omterm{def:b3:topology:compact}{compact}, covered by
finitely many others) and \omterm{def:b3:topology:hausdorff}{Hausdorff} (separate $x$ from $\infty$
by a \omterm{def:b3:topology:compact}{compact} \omterm{def:b3:topology:topology}{neighborhood} of $x$ and its complement). Example:
$\hat\R \cong S^1$, and $\widehat{\R^n} \cong S^n$ by
stereographic projection (\cref{exo:b3:topology:11}).
\end{definition}

\section{Connectedness}

\begin{definition}\label{def:b3:topology:connected}
$X$ is \emph{connected}\index{connected space} if it is not the
union of two disjoint nonempty \omterm{def:b3:topology:topology}{open sets} --- equivalently, its
only subsets both open and closed are $\varnothing$ and $X$;
equivalently, every \omterm{def:b3:topology:continuity}{continuous map} $X \to \{0, 1\}$ (discrete)
is constant. A subset is connected if it is as a subspace.
\end{definition}

\begin{theorem}\label{thm:b3:topology:connectedbasics}
\begin{enumerate}
  \item The \omterm{def:b3:topology:connected}{connected} subsets of $\R$ are exactly the intervals.
  \item \omterm{def:b3:topology:continuity}{Continuous} images of \omterm{def:b3:topology:connected}{connected} sets are \omterm{def:b3:topology:connected}{connected} (whence
        the intermediate value theorem: a \omterm{def:b3:topology:continuity}{continuous} real map on
        a \omterm{def:b3:topology:connected}{connected space} has an interval as image).
  \item If $(A_i)$ are \omterm{def:b3:topology:connected}{connected} with a common point, $\bigcup
        A_i$ is \omterm{def:b3:topology:connected}{connected}. If $A$ is \omterm{def:b3:topology:connected}{connected} and $A \subseteq
        B \subseteq \bar A$, then $B$ is \omterm{def:b3:topology:connected}{connected}.
  \item Finite products of \omterm{def:b3:topology:connected}{connected spaces} are \omterm{def:b3:topology:connected}{connected}.
\end{enumerate}
\end{theorem}

\begin{proof}
Throughout we use the $\{0,1\}$-criterion: $f \colon X \to
\{0,1\}$ \omterm{def:b3:topology:continuity}{continuous} must be constant.

(1) A non-interval $A$ misses some $c$ between $a, b \in A$:
$A = (A \cap (-\infty, c)) \sqcup (A \cap (c, +\infty))$
disconnects. Conversely let $I$ be an interval and $f \colon I
\to \{0,1\}$ \omterm{def:b3:topology:continuity}{continuous} with $f(a) = 0$, $f(b) = 1$, $a < b$.
Let $c = \sup\{x \in [a,b] : f(x) = 0\}$; \omterm{def:b3:topology:continuity}{continuity} at $c$
forces $f(c) = 0$ (limit of values $0$: every \omterm{def:b3:topology:topology}{neighborhood} of
$c$ meets $\{f = 0\}$; $\{f=0\}$ is closed) and then $c < b$
with $f \equiv 1$ on $(c, b]$, so $f(c) = 1$ by the same
\omterm{def:b3:topology:interior}{closure} argument on $\{f = 1\} \ni c$: contradiction.

(2) A \omterm{def:b3:topology:continuity}{continuous} $g \colon f(X) \to \{0,1\}$ yields $g \circ f$
constant, so $g$ is constant on the image.

(3) A \omterm{def:b3:topology:continuity}{continuous} $f\colon \bigcup A_i \to \{0,1\}$ is constant
on each $A_i$, with the same value at the common point. For the
\omterm{def:b3:topology:interior}{closure}: $f \colon B \to \{0,1\}$ is constant $= c$ on $A$; any
$x \in B \subseteq \bar A$ is in the \omterm{def:b3:topology:interior}{closure} of $A$, and $f^{-1}
(f(x))$ is a \omterm{def:b3:topology:topology}{neighborhood} of $x$ (preimage of open), which must
meet $A$: $f(x) = c$.

(4) For $X \times Y$ and $f\colon X\times Y \to \{0,1\}$: any
two points $(x,y), (x',y')$ are joined by the ``elbow''
$\{x\}\times Y \cup X \times \{y'\}$, a union of two \omterm{def:b3:topology:connected}{connected}
sets (\omterm{def:b3:topology:continuity}{homeomorphic} to $Y$, $X$) meeting at $(x, y')$: by (3) and
(2), $f$ agrees on the two points.
\end{proof}

\begin{definition}\label{def:b3:topology:pathconnected}
$X$ is \emph{path-connected}\index{path-connected space} if any
two points are joined by a \emph{path} (\omterm{def:b3:topology:continuity}{continuous} $\gamma
\colon [0,1] \to X$). \omterm{def:b3:topology:connected}{Path-connected} implies \omterm{def:b3:topology:connected}{connected}: two
values of a \omterm{def:b3:topology:continuity}{continuous} $f \colon X \to \{0,1\}$ at $x, y$ are
values of the constant $f\circ\gamma$
(\cref{thm:b3:topology:connectedbasics}(1)--(2)). Convex subsets
of normed spaces are \omterm{def:b3:topology:connected}{path-connected} (segments); so are $\R^n
\setminus \{0\}$ for $n \geq 2$ (go around the origin), and
$S^n$ for $n \geq 1$ (project paths from $\R^{n+1}\setminus
\{0\}$).
\end{definition}

\begin{example}[The topologist's sine curve]
\label{ex:b3:topology:sinecurve}
Let $\Gamma = \{(x, \sin\frac1x) : 0 < x \leq 1\}$ and $S =
\bar\Gamma = \Gamma \cup (\{0\}\times[-1,1])$ (every point $(0,
y)$, $\abs y \leq 1$, is a limit of points of $\Gamma$: solve
$\sin\frac1x = y$ near $0$). Then $S$ is
\emph{\omterm{def:b3:topology:connected}{connected}} --- \omterm{def:b3:topology:interior}{closure} of the \omterm{def:b3:topology:connected}{connected} $\Gamma$, a
\omterm{def:b3:topology:continuity}{continuous} image of $(0, 1]$
(\cref{thm:b3:topology:connectedbasics}(3)) --- but \emph{not
\omterm{def:b3:topology:pathconnected}{path-connected}}: a \omterm{def:b3:topology:pathconnected}{path} from $(1, \sin 1)$ to $(0,0)$ would have
to traverse abscissas $\to 0$ while the ordinate oscillates
between $\pm1$; \cref{exo:b3:topology:9} makes this rigorous.
\omterm{def:b3:topology:connected}{Connectedness} and \omterm{def:b3:topology:connected}{path-connectedness} genuinely
differ.\index{topologist's sine curve}
\end{example}

\begin{omfigure}
\begin{tikzpicture}
\begin{axis}[omaxis, width=10.5cm, height=4.6cm,
    xmin=-0.02, xmax=0.55, ymin=-1.25, ymax=1.25,
    xtick={0, 0.25, 0.5}, ytick={-1, 1}]
  \addplot[omThm, thick, domain=0.006:0.55, samples=1200]
    {sin(deg(1/x))};
  \addplot[omDef, very thick] coordinates {(0,-1) (0,1)};
\end{axis}
\end{tikzpicture}

{\small The \omterm{ex:b3:topology:sinecurve}{topologist's sine curve}: the graph of
$\sin\frac1x$ (red) accumulates on the whole segment
$\{0\}\times[-1,1]$ (blue). The union is \omterm{def:b3:topology:connected}{connected} but not
\omterm{def:b3:topology:pathconnected}{path-connected}: no \omterm{def:b3:topology:continuity}{continuous} \omterm{def:b3:topology:pathconnected}{path} can cross the infinitely
many oscillations in finite parameter time.}
\end{omfigure}

\begin{definition}\label{def:b3:topology:components}
The \emph{connected component}\index{connected component} of $x
\in X$ is the union of all \omterm{def:b3:topology:connected}{connected} subsets containing $x$ ---
the largest one (\cref{thm:b3:topology:connectedbasics}(3)).
Components partition $X$ and are closed (\omterm{def:b3:topology:interior}{closures} of \omterm{def:b3:topology:connected}{connected}
sets are \omterm{def:b3:topology:connected}{connected}). $X$ is \emph{totally disconnected} if all
components are singletons ($\Q$; the Cantor set of the weekend
problem).
\end{definition}

\begin{theorem}\label{thm:b3:topology:rnotr2}
$\R$ is \omterm{def:b3:topology:continuity}{homeomorphic} to no $\R^n$ with $n \geq 2$.
\end{theorem}

\begin{proof}
Suppose $h \colon \R \to \R^n$ is a \omterm{def:b3:topology:continuity}{homeomorphism}. Removing a
point: $\R \setminus \{0\}$ is \omterm{def:b3:topology:continuity}{homeomorphic} to $\R^n \setminus
\{h(0)\}$. But $\R\setminus\{0\}$ is disconnected, while
$\R^n\setminus\{\text{point}\}$ is \omterm{def:b3:topology:pathconnected}{path-connected} for $n \geq
2$ (\cref{def:b3:topology:pathconnected}; translate the point to
$0$): \omterm{def:b3:topology:connected}{connectedness} is a \omterm{def:b3:topology:continuity}{homeomorphism} invariant
(\cref{thm:b3:topology:connectedbasics}(2)) --- contradiction.
(That $\R^2 \not\cong \R^3$ requires finer invariants ---
\omterm{def:b3:galois:algebraic}{algebraic} \omterm{def:b3:topology:topology}{topology}; the weekend problem shows the danger:
\emph{\omterm{def:b3:topology:continuity}{continuous} surjections} $\R \to \R^2$ do exist.)
\end{proof}

\begin{method}\label{met:b3:topology:connectargument}
To prove a set is \omterm{def:b3:topology:connected}{connected}: exhibit it as a \omterm{def:b3:topology:continuity}{continuous} image,
a union of overlapping \omterm{def:b3:topology:connected}{connected} sets, a \omterm{def:b3:topology:interior}{closure}, or a product
(\cref{thm:b3:topology:connectedbasics}); for subsets of normed
spaces, prove \omterm{def:b3:topology:connected}{path-connectedness} with explicit \omterm{def:b3:topology:pathconnected}{paths} (segments,
arcs, elbows). To prove two spaces are \emph{not} \omterm{def:b3:topology:continuity}{homeomorphic}:
find a \omterm{def:b3:topology:topology}{topological} invariant that differs --- compactness,
\omterm{def:b3:topology:connected}{connectedness}, number of components, or components after
deleting a well-chosen finite set (the trick of
\cref{thm:b3:topology:rnotr2}: it also proves $[0,1) \not\cong
(0,1)$ and $S^1 \not\cong [0,1]$).
\end{method}

\section{Exercises}

\begin{exercise}[$\star$]\label{exo:b3:topology:1}
(a) List all topologies on $\{a, b\}$, classify them up to
\omterm{def:b3:topology:continuity}{homeomorphism}, and determine which ones are \omterm{def:b3:topology:connected}{connected} and which
are \omterm{def:b3:topology:hausdorff}{Hausdorff}.
(b) In $\R$ (usual \omterm{def:b3:topology:topology}{topology}), compute \omterm{def:b3:topology:interior}{interior}, \omterm{def:b3:topology:interior}{closure} and
boundary of $\Q$, of $[0,1) \cup \{2\}$, and of $\{1/n : n \geq
1\}$.
\end{exercise}

\begin{exercise}[$\star$]\label{exo:b3:topology:2}
(a) Show that $f\colon X \to Y$ is \omterm{def:b3:topology:continuity}{continuous} iff $f^{-1}(B)$
is open for every $B$ in a fixed \omterm{def:b3:topology:basis}{basis} of $Y$.
(b) Show that $f = \mathbf 1_{[0,\infty)} \colon \R \to \R$ is
not \omterm{def:b3:topology:continuity}{continuous} but is \emph{\omterm{def:b3:topology:continuity}{right-continuous}}. Verify that the
half-open intervals $[a, b)$ form a \omterm{def:b3:topology:basis}{basis of a topology} on the
source $\R$ (the \emph{Sorgenfrey line}), that this \omterm{def:b3:topology:topology}{topology} is
strictly finer than the usual one, and that a function $\R \to
\R$ (usual target) is \omterm{def:b3:topology:continuity}{continuous} from the Sorgenfrey line iff it
is \omterm{def:b3:topology:continuity}{right-continuous} at every point.
\end{exercise}

\begin{exercise}[$\star$]\label{exo:b3:topology:3}
On an infinite set with the cofinite \omterm{def:b3:topology:topology}{topology}, show: any two
nonempty \omterm{def:b3:topology:topology}{open sets} meet (so the space is not \omterm{def:b3:topology:hausdorff}{Hausdorff}, hence
not metrizable); every injective sequence converges to
\emph{every} point. Where does the uniqueness-of-limits proof
use \omterm{def:b3:topology:hausdorff}{Hausdorff}?
\end{exercise}

\begin{exercise}[$\star\star$]\label{exo:b3:topology:4}
(a) Show that the projections $\pi_X, \pi_Y$ of a product are
\omterm{def:b3:topology:continuity}{continuous} and \emph{open} (images of opens are open), but not
closed in general ($\{xy = 1\}$ in $\R^2$).
(b) Show that a sequence in a countable product $\prod_n X_n$ of
metric spaces converges iff each coordinate converges, and that
$d(x, y) = \sum_n 2^{-n}\min(1, d_n(x_n, y_n))$ metrizes the
\omterm{def:b3:topology:constructions}{product topology}.
(c) On $\R^\N$, show that the ``box \omterm{def:b3:topology:topology}{topology}'' (all products of
opens are open) is strictly finer: the sequence
$x^{(k)} = (1/k, 1/k, \dots)$ converges to $0$ in the \omterm{def:b3:topology:constructions}{product
topology} but not in the box \omterm{def:b3:topology:topology}{topology}.
\end{exercise}

\begin{exercise}[$\star\star$]\label{exo:b3:topology:5}
The circle, three ways. Show that the following are pairwise
\omterm{def:b3:topology:continuity}{homeomorphic}, with explicit maps: (i) $S^1 \subseteq \R^2$;
(ii) $\R/\Z$ (\omterm{def:b3:topology:constructions}{quotient topology}); (iii) $[0,1]/(0 \sim 1)$.
\emph{(For (ii): show $\R/\Z$ is \omterm{def:b3:topology:hausdorff}{Hausdorff} --- lift two classes
to representatives at distance $\leq \frac12$ --- and that
$\pi([0,1])$ is everything, so the quotient is \omterm{def:b3:topology:compact}{compact}; then use
\cref{cor:b3:topology:compacthomeo}.)}
\end{exercise}

\begin{exercise}[$\star\star$]\label{exo:b3:topology:6}
Let $X$ be a metric space.
(a) Show that a finite union of \omterm{def:b3:topology:compact}{compact} subsets is \omterm{def:b3:topology:compact}{compact}, and
that an arbitrary intersection is.
(b) If $K$ is \omterm{def:b3:topology:compact}{compact}, $F$ closed, $K \cap F = \varnothing$,
show $d(K, F) = \inf\{d(x,y) : x \in K, y \in F\} > 0$; give a
counterexample with two disjoint closed sets.
(c) Show that $X$ is \omterm{def:b3:topology:compact}{compact} iff \emph{every} \omterm{def:b3:topology:continuity}{continuous} $f
\colon X \to \R$ is bounded. \emph{(If some sequence has no
convergent subsequence, build an unbounded \omterm{def:b3:topology:continuity}{continuous} function
supported near its terms; or use $x \mapsto
d(x, \cdot)$-type functions --- one clean route: if $(x_n)$ has
no cluster point, the set $\{x_n\}$ is closed and discrete, and
$f(x_n) = n$ extends \omterm{def:b3:topology:continuity}{continuously} by Tietze-free means:
$f(x) = \sum_n n\,\max\bigl(0, 1 - \frac{d(x, x_n)}{r_n}\bigr)$
for small enough $r_n$.)}
\end{exercise}

\begin{exercise}[$\star\star$]\label{exo:b3:topology:7}
(Lebesgue number) Let $(U_i)$ be an open cover of a \omterm{def:b3:topology:compact}{compact}
metric space $X$. Show there is $\delta > 0$ such that every
subset of diameter $< \delta$ lies in a single $U_i$.
\emph{(Otherwise pick $A_n$ of diameter $< 1/n$ in no $U_i$, and
a cluster point of chosen $x_n \in A_n$.)} Deduce Heine's
theorem (\cref{cor:b3:topology:heineborel}(b)) again.
\end{exercise}

\begin{exercise}[$\star\star$]\label{exo:b3:topology:8}
(a) Show that $GL_n(\R)$ is an open, \omterm{def:b3:topology:interior}{dense} subset of $M_n(\R)$,
and that it is disconnected: the sign of the determinant
separates it into (at least) two pieces. Show on the other hand
that $GL_n(\C)$ is \omterm{def:b3:topology:pathconnected}{path-connected}.
\emph{(For $A, B \in GL_n(\C)$: $\lambda \mapsto \det\bigl((1 -
\lambda)A + \lambda B\bigr)$ is a polynomial, not identically
zero, so it has finitely many roots in $\C$: pick a \omterm{def:b3:topology:pathconnected}{path} of
$\lambda$'s in $\C$ from $0$ to $1$ avoiding them.)}
(b) Prove density: $A + \varepsilon I$ is invertible for small
$\varepsilon > 0$.
\end{exercise}

\begin{exercise}[$\star\star$]\label{exo:b3:topology:9}
(a) Show that the \omterm{def:b3:topology:components}{connected components} of $\Q$ are the
singletons, and that $\Q$ is not \omterm{def:b3:topology:locallycompact}{locally compact} (a \omterm{def:b3:topology:compact}{compact}
\omterm{def:b3:topology:topology}{neighborhood} of $0$ in $\Q$ would contain $[-r, r]\cap\Q$,
whose \omterm{def:b3:topology:interior}{closure} in $\Q$ is not \omterm{def:b3:topology:compact}{compact}: cut at an irrational).
(b) Complete \cref{ex:b3:topology:sinecurve}: no \omterm{def:b3:topology:pathconnected}{path} joins
$(0,0)$ to $\Gamma$. \emph{(If $\gamma = (\gamma_1, \gamma_2)$
is such a \omterm{def:b3:topology:pathconnected}{path}, let $t^* = \sup\{t : \gamma_1(t) = 0\}$; for $t
> t^*$ the point moves on the graph; choose $t_k \downarrow
t^*$ with $\gamma_1(t_k) \to 0$ hitting abscissas where $\sin$
is alternately $\pm1$ --- the intermediate value theorem
supplies them --- contradicting \omterm{def:b3:topology:continuity}{continuity} of $\gamma_2$ at
$t^*$.)}
\end{exercise}

\begin{exercise}[$\star\star\star$]\label{exo:b3:topology:10}
The middle-thirds Cantor set $C = \bigcap_n C_n$, where $C_0 =
[0,1]$ and $C_{n+1}$ removes the open middle third of each
interval of $C_n$.
(a) Show $C = \bigl\{\sum_{n\geq1} a_n3^{-n} : a_n \in
\{0,2\}\bigr\}$, and that $C$ is \omterm{def:b3:topology:compact}{compact}, with empty \omterm{def:b3:topology:interior}{interior},
and has no isolated point (\emph{\omterm{prop:b3:galois:perfect}{perfect}}).
(b) Show that $C$ is totally disconnected.
(c) Show that $x \mapsto (a_n(x))_n$ is a \omterm{def:b3:topology:continuity}{homeomorphism} $C \to
\{0,2\}^\N$ (\omterm{def:b3:topology:constructions}{product topology} on the discrete two-point
space); deduce that $C$ is uncountable, and that $C \times C
\cong C$.
\end{exercise}

\begin{exercise}[$\star\star\star$]\label{exo:b3:topology:11}
Stereographic projection: from the north pole $N = (0, \dots,
0, 1)$ of $S^n \subseteq \R^{n+1}$, the map $\sigma(x) =
\frac{(x_1, \dots, x_n)}{1 - x_{n+1}}$ is a \omterm{def:b3:topology:continuity}{homeomorphism} $S^n
\setminus \{N\} \to \R^n$ (give the inverse explicitly). Deduce
that $S^n$ is \omterm{def:b3:topology:continuity}{homeomorphic} to the \omterm{def:b3:topology:locallycompact}{one-point compactification}
$\widehat{\R^n}$, and that removing \emph{any} point of $S^n$
leaves a space \omterm{def:b3:topology:continuity}{homeomorphic} to $\R^n$.
\end{exercise}

\begin{exercise}[$\star\star\star$]\label{exo:b3:topology:12}
(The \omterm{ex:b3:topology:sinecurve}{topologist's sine curve}) Let
\[
T = \bigl\{(x, \sin\tfrac1x) : 0 < x \leq 1\bigr\}
\cup \bigl(\{0\}\times\intcc{-1}1\bigr) \subseteq \R^2 .
\]
(a) Show that $T$ is \omterm{def:b3:topology:compact}{compact}, and that it is the \omterm{def:b3:topology:interior}{closure} of
the graph part.
(b) Show that $T$ is \omterm{def:b3:topology:connected}{connected} \emph{(the graph is \omterm{def:b3:topology:connected}{connected}
as a \omterm{def:b3:topology:continuity}{continuous} image; its \omterm{def:b3:topology:interior}{closure} remains \omterm{def:b3:topology:connected}{connected})}.
(c) Show that $T$ is \emph{not} \omterm{def:b3:topology:pathconnected}{path-connected}: no \omterm{def:b3:topology:continuity}{continuous}
\omterm{def:b3:topology:pathconnected}{path} joins $(0,0)$ to $(1, \sin1)$. \emph{(If $\gamma =
(\gamma_1, \gamma_2)$ is such a \omterm{def:b3:topology:pathconnected}{path}, let $t^* = \sup\{t :
\gamma_1(t) = 0\}$; just after $t^*$, $\gamma_1$ takes all
small positive values (intermediate value theorem), so
$\gamma_2$ oscillates between $\pm1$ on every interval
$(t^*, t^* + \delta)$ --- contradict \omterm{def:b3:topology:continuity}{continuity} at $t^*$.)}
(d) Conclude that \omterm{def:b3:topology:connected}{path-connectedness} is strictly stronger
than \omterm{def:b3:topology:connected}{connectedness}, and show that no such example can be
open in $\R^n$: an \emph{open} \omterm{def:b3:topology:connected}{connected} subset of $\R^n$ is
\omterm{def:b3:topology:pathconnected}{path-connected} \emph{(the set of points joinable to a base
point by a \omterm{def:b3:topology:pathconnected}{path} is open and closed in the domain)}.
\end{exercise}

\section{Problem: the Cantor set and a space-filling curve}

\begin{problem}[{Weekend problem --- Peano curves exist, and why
they are not homeomorphisms}]\label{pb:b3:topology:1}
In 1890 Peano stunned analysis with a \emph{\omterm{def:b3:topology:continuity}{continuous}
surjection} $[0,1] \to [0,1]^2$: a curve filling a square. We
build one with bare hands from the Cantor set $C$ of
\cref{exo:b3:topology:10} (whose results may be used freely),
and then prove that no such map can be injective: squares are
not curves. Throughout, elements of $C$ are written $x =
\sum_{n \geq 1} \frac{2b_n(x)}{3^n}$ with digits $b_n(x) \in
\{0, 1\}$.

\medskip
\noindent\textbf{Part I --- Reading off digits \omterm{def:b3:topology:continuity}{continuously}.}
\begin{enumerate}
  \item Show that if $x, y \in C$ satisfy $\abs{x - y} <
        3^{-m}$, then $b_n(x) = b_n(y)$ for all $n \leq m$.
        \emph{(If the first differing digit is $b_k$, then
        $\abs{x - y} \geq \frac{2}{3^k} - \sum_{j>k}\frac2{3^j}
        = \frac1{3^k}$.)}
  \item Deduce that each digit function $b_n \colon C \to \{0,
        1\}$ is \omterm{def:b3:topology:continuity}{continuous}, and re-derive the \omterm{def:b3:topology:continuity}{homeomorphism} $C
        \cong \{0,1\}^\N$ of \cref{exo:b3:topology:10}(c).
\end{enumerate}

\medskip
\noindent\textbf{Part II --- A \omterm{def:b3:topology:continuity}{continuous} surjection $C \to
[0,1]^2$.}
\begin{enumerate}[resume]
  \item Define $g \colon C \to [0,1]$ by $g(x) = \sum_{n\geq1}
        b_n(x)\,2^{-n}$ (read the Cantor digits as binary).
        Show that $g$ is \omterm{def:b3:topology:continuity}{continuous} (use question 1) and
        surjective. Is it injective?
  \item Define $\Phi \colon C \to [0,1]^2$ by
        \[
        \Phi(x) = \Bigl(\ \sum_{n\geq1} b_{2n-1}(x)\,2^{-n},\
        \sum_{n\geq1} b_{2n}(x)\,2^{-n}\ \Bigr)
        \]
        (odd digits give the abscissa, even digits the
        ordinate). Show that $\Phi$ is \omterm{def:b3:topology:continuity}{continuous} and
        surjective.
\end{enumerate}

\medskip
\noindent\textbf{Part III --- Filling the gaps: the Peano
curve.}
\begin{enumerate}[resume]
  \item The complement $[0,1] \setminus C$ is a countable union
        of disjoint open intervals $(u, v)$ (the removed middle
        thirds) whose endpoints lie in $C$. Define $f \colon
        [0,1] \to [0,1]^2$ by $f = \Phi$ on $C$, extended
        \emph{affinely} on each gap:
        \[
        f(t) = \frac{v - t}{v - u}\,\Phi(u) +
        \frac{t - u}{v - u}\,\Phi(v),
        \qquad t \in (u, v).
        \]
        Show that $f$ is well defined and surjective onto
        $[0,1]^2$.
  \item Show that $f$ is \omterm{def:b3:topology:continuity}{continuous} at every point of $[0,1]
        \setminus C$ (locally affine), and at every point of
        $C$: given $\varepsilon$, pick $m$ with $2^{-m} <
        \varepsilon/2$ and use question 1 to control $\Phi$ on
        $C$ near $x$, then check that the affine
        interpolation cannot escape: on a gap $(u,v) \subseteq
        (x - 3^{-m}, x + 3^{-m})$, the values $f(t)$ lie on the
        segment $[\Phi(u), \Phi(v)]$, both ends close to
        $\Phi(x)$. Conclude: $f$ is a \emph{\omterm{def:b3:topology:continuity}{continuous}
        surjection} $[0,1] \to [0,1]^2$.
  \item Deduce \omterm{def:b3:topology:continuity}{continuous} surjections $[0,1] \to [0,1]^n$ for
        every $n \geq 2$, and $\R \to \R^2$.
\end{enumerate}

\medskip
\noindent\textbf{Part IV --- But never injective.}
\begin{enumerate}[resume]
  \item Show that a \omterm{def:b3:topology:continuity}{continuous} \emph{injection} $f \colon [0,1]
        \to [0,1]^2$ would be a \omterm{def:b3:topology:continuity}{homeomorphism} onto its image
        (\cref{cor:b3:topology:compacthomeo}).
  \item Show that $[0,1]$ and $[0,1]^2$ are not \omterm{def:b3:topology:continuity}{homeomorphic}:
        remove a well-chosen point and compare \omterm{def:b3:topology:connected}{connectedness}
        (\cref{met:b3:topology:connectargument}).
  \item Conclude: a \omterm{def:b3:topology:continuity}{continuous} surjection $[0,1] \to [0,1]^2$
        can never be injective --- an injective one would make
        $[0,1]^2 = f([0,1])$ \omterm{def:b3:topology:continuity}{homeomorphic} to $[0,1]$ by
        question 8, contradicting question 9. Where exactly did
        compactness of $[0,1]$ enter the argument?
  \item (Culmination) Assemble the moral: there is a \omterm{def:b3:topology:continuity}{continuous}
        surjection but no \omterm{def:b3:topology:continuity}{continuous} bijection $[0,1] \to
        [0,1]^2$. What does this say about ``dimension'' as a
        \omterm{def:b3:topology:topology}{topological} notion? Formulate precisely one theorem
        proved in this problem and one plausible statement that
        remains beyond our tools (invariance of domain).
\end{enumerate}

\medskip
\noindent\textbf{Part V --- The arithmetic of the Cantor set.}
\begin{enumerate}[resume]
  \item Prove that $C + C = [0, 2]$: given $t \in [0, 1]$,
        write $t = \sum_n d_n3^{-n}$ with digits $d_n \in \{0,
        1, 2\}$ and split each digit as $d_n = a_n + b_n$ with
        $a_n, b_n \in \{0, 1\}$; conclude that $\frac C2 +
        \frac C2 = [0, 1]$ \emph{(which subset of $[0,1]$ is
        $\frac C2$, in terms of ternary digits?)}, then rescale.
        No carries are ever needed --- say why this is the crux.
  \item Interpret geometrically: the ``Cantor dust'' $C \times
        C \subseteq \R^2$, of empty \omterm{def:b3:topology:interior}{interior}, casts a
        \emph{full} shadow on the diagonal: the projection $(x,
        y) \mapsto x + y$ maps it onto $[0, 2]$. Also record
        the self-similarity $C = \frac C3 \cup \bigl(\frac C3 +
        \frac23\bigr)$, the equation behind every picture of
        $C$.
  \item Show that digit interleaving defines a \omterm{def:b3:topology:continuity}{homeomorphism} $C
        \to C \times C$, and deduce $C \cong C^n$ for every $n$
        --- the Cantor set is its own square, cube, \dots{}
        Which familiar spaces share this property?
  \item Show that $\frac14 \in C$ \emph{(compute its ternary
        expansion: $\frac14 = \sum_{k\geq1} 2\cdot3^{-2k}$)}
        although $\frac14$ is an endpoint of no removed
        interval; deduce --- counting endpoints --- that
        endpoints form a countable, proper subset of $C$.
  \item Show that the binary-reading map $g$ of question 3 is
        at most $2$-to-$1$, and describe exactly which points
        of $[0,1]$ have two preimages. ($g$ collapses $C$ onto
        $[0,1]$ by gluing countably many pairs: the
        combinatorial shadow of the Cantor staircase met again
        in \cref{ch:b3:measure}.)
\end{enumerate}

\medskip
\noindent\textbf{Part VI --- \omterm{prop:b3:galois:perfect}{Perfect} \omterm{def:b3:topology:compact}{compact} sets: Cantor
everywhere.} A nonempty \omterm{def:b3:topology:compact}{compact} metric space $P$ is
\emph{\omterm{prop:b3:galois:perfect}{perfect}} if it has no isolated point.
\begin{enumerate}[resume]
  \item Let $P$ be \omterm{prop:b3:galois:perfect}{perfect} \omterm{def:b3:topology:compact}{compact} and let $x \in P$, $r > 0$.
        Show that the ball $B(x, r)$ contains two points $y_0
        \neq y_1$ of $P$, and hence two \emph{disjoint} closed
        balls $\bar B(y_0, \rho)$, $\bar B(y_1, \rho)$ inside
        $B(x, r)$, centered at points of $P$, of radius $\rho
        \leq r/4$ as small as desired. Explain why perfection
        (no isolated points) is exactly what allows this
        splitting to be repeated inside \emph{each} of the two
        new balls.
  \item Iterate: build closed sets $F_w \neq \varnothing$
        indexed by finite binary words $w$, with $F_{w0},
        F_{w1} \subseteq F_w$ disjoint and
        $\operatorname{diam}F_w \leq 2^{-\abs w}
        \operatorname{diam}P$. Show that for every infinite
        word $\varepsilon \in \{0,1\}^\N$ the intersection
        $\bigcap_mF_{\varepsilon_1\cdots\varepsilon_m}$ is a
        single point $h(\varepsilon)$ \emph{(finite
        intersection property of the \omterm{def:b3:topology:compact}{compact} $P$)}.
  \item Show that $h\colon\{0,1\}^\N \to P$ is injective and
        \omterm{def:b3:topology:continuity}{continuous}, and conclude: \emph{every \omterm{prop:b3:galois:perfect}{perfect} \omterm{def:b3:topology:compact}{compact}
        metric space is uncountable} --- in fact of cardinality
        at least that of $\{0,1\}^\N$. Recover: $[0,1]$ and $C$
        are uncountable.
  \item Show that $h$ is a \omterm{def:b3:topology:continuity}{homeomorphism} onto its image
        \emph{(\omterm{def:b3:topology:continuity}{continuous} injection from a \omterm{def:b3:topology:compact}{compact},
        \cref{cor:b3:topology:compacthomeo})}: every \omterm{prop:b3:galois:perfect}{perfect}
        \omterm{def:b3:topology:compact}{compact} metric space contains a \omterm{def:b3:topology:continuity}{homeomorphic} copy of
        the Cantor set. The Cantor set is not an oddity but the
        universal germ of \omterm{def:b3:topology:compact}{compact} perfection.
  \item Deduce that every countable \omterm{def:b3:topology:compact}{compact} metric space has an
        isolated point, and exhibit one where the isolated
        points are \omterm{def:b3:topology:interior}{dense} but not everything: $K = \{0\} \cup
        \{\frac1n : n \geq 1\}$.
  \item (Cantor--Bendixson for $\R$) Let $F \subseteq \R$ be
        closed. Call $x$ a \emph{condensation point} of $F$ if
        every \omterm{def:b3:topology:topology}{neighborhood} of $x$ meets $F$ uncountably. Show
        that the condensation points of an uncountable closed
        $F$ form a nonempty \omterm{prop:b3:galois:perfect}{perfect} closed set $P$, and $F
        \setminus P$ is countable \emph{(cover the
        non-condensation points by countably many rational
        intervals meeting $F$ countably)}. Conclude: every
        closed subset of $\R$ is countable or has the
        cardinality of the continuum --- the continuum
        hypothesis holds for closed sets.
\end{enumerate}

\medskip
\noindent\textbf{Part VII --- Codas: how regular, and how
far from \omterm{prop:b3:galois:perfect}{perfect}.}
\begin{enumerate}[resume]
  \item (H\"older regularity of the curve) Set $\alpha =
        \frac{\ln 2}{2\ln 3} \approx 0.3155$, and note the
        identity $3^\alpha = \sqrt2$. Using question 1, show
        that
        \[
        \norm{\Phi(x) - \Phi(y)}_\infty \leq 2\,\abs{x -
        y}^\alpha \qquad (x, y \in C),
        \]
        then propagate the estimate through the affine gaps:
        show that the Peano curve $f$ of question 6 satisfies
        $\norm{f(t) - f(s)}_\infty \leq 6\abs{t - s}^\alpha$ on
        all of $[0,1]$ \emph{(treat a pair in the same gap by
        interpolation, then a general pair by passing through
        the extreme points of $C \cap [t, s]$)}.
  \item (The exponent $\frac12$ is a wall) Show that no
        surjection $F \colon [0,1] \to [0,1]^2$ can be
        $\alpha$-H\"older with $\alpha > \frac12$: cut $[0,1]$
        into $N$ intervals, bound the diameters of their
        images, and count the points of the grid
        $\bigl(\frac iM, \frac jM\bigr)$, $0 \leq i, j \leq M$,
        that a set of diameter $< \frac1M$ can contain; choose
        $M$ of order $N^\alpha$ and let $N \to \infty$. Locate
        our curve ($\alpha \approx 0.3155$) relative to the
        wall, and record without proof that Hilbert's curve
        achieves the critical exponent $\frac12$.
  \item (Derived sets: measuring imperfection) For $F$ closed
        in $\R$, let $F'$ be the set of limit points of $F$ (a
        closed subset), and iterate: $F^{(0)} = F$, $F^{(i+1)}
        = (F^{(i)})'$. Verify that
        \[
        K = \{0\} \cup \Bigl\{\frac1m\Bigr\}_{m \geq 1} \cup
        \Bigl\{\frac1m + \frac1n : m \geq 1,\ n > m^2\Bigr\}
        \]
        is a countable \omterm{def:b3:topology:compact}{compact} set with $K' = \{0\} \cup
        \{\frac1m\}$, $K'' = \{0\}$, $K''' = \varnothing$
        \emph{(check that the $m$-th cluster lives in the
        interval $(\frac1m, \frac1{m-1})$, so the clusters do
        not interleave)}, and that its isolated points are
        \omterm{def:b3:topology:interior}{dense} in $K$, as question 21 predicts. Describe the
        induction producing, for every $j$, a countable \omterm{def:b3:topology:compact}{compact}
        $K_j$ with $K_j^{(j)} = \{0\}$ and $K_j^{(j+1)} =
        \varnothing$, and contrast with the \omterm{prop:b3:galois:perfect}{perfect} sets of
        question 22, for which the derivation never moves:
        finite rank is the exact opposite of perfection.
\end{enumerate}
\end{problem}
